\begin{figure}[htbp]
\centering
\begin{tikzpicture}[x=1cm,y=1cm]
\node[dissertation block,text width=2.85cm] (joint) at (-5.55,0.35)
  {Представлення\\$Z$};
\node[dissertation block,text width=2.85cm] (decision) at (-2.15,0.35)
  {Функція рішення\\$F_D$};
\node[dissertation block,text width=2.85cm] (state) at (1.25,0.35)
  {Аналітичний стан\\$S_A$};
\node[dissertation block,text width=3.40cm] (classification) at (5.05,2.15)
  {Визначення класу\\$p,\hat y,\hat k$};
\node[dissertation block,text width=3.40cm] (risk) at (5.05,0.35)
  {Ризик і упевненість\\$\rho_A,c_A$};
\node[dissertation block,text width=3.40cm] (explanation) at (5.05,-1.45)
  {Перевірюване пояснення\\$\Pi$};

\draw[dissertation arrow] (joint.east) -- (decision.west);
\draw[dissertation arrow] (decision.east) -- (state.west);
\draw[dissertation arrow] (state.east) -- ++(0.45,0) |- (classification.west);
\draw[dissertation arrow] (state.east) -- (risk.west);
\draw[dissertation arrow] (state.east) -- ++(0.45,0) |- (explanation.west);
\end{tikzpicture}
\caption{Формування аналітичного стану моделі без підміни його технологічною дією}
\label{fig:decision_criteria}
\end{figure}
